\documentclass[11pt]{beamer}
\usetheme{Madrid}
\usepackage[T1]{fontenc}
\usepackage[utf8]{inputenc}
\usepackage{amsmath,amsfonts,amssymb}
\usepackage{graphicx}
\usepackage{xcolor}
\usepackage{tikz}
\usepackage{array}
\usepackage{longtable}
\usepackage{booktabs}
\usepackage{hyperref}

\title{EB2-NIW \& EB1A Profile Building}
\subtitle{High-Volume Research Strategy: 20-50 Papers  1000-2000 Citations 2 years Under 20k}
\author{Complete DIY Research Profile System}
\date{\today}

\setbeamercovered{transparent}
\setbeamertemplate{navigation symbols}{}
\setbeamerfont{itemize/enumerate body}{size=\scriptsize}
\setbeamerfont{itemize/enumerate subbody}{size=\tiny}
\setbeamerfont{block body}{size=\scriptsize}
\setbeamerfont{block title}{size=\normalsize}

\begin{document}
	
	% ============================================================
	% SLIDE 1: TITLE
	% ============================================================
	\begin{frame}
		\titlepage
		\centering
		\scriptsize \textit{20-50 Papers  1000-2000 Citations   \$15-25K Budget   Government Impact Focus}
	\end{frame}
	
	% ============================================================
	% SLIDE 2: THE TARGET NUMBERS
	% ============================================================
	\begin{frame}{The Target Numbers - What You Need}
		\scriptsize
		\begin{block}{Publication Targets}
			\begin{itemize}
				\item \textbf{20-50 total publications} (journals + conferences + book chapters)
				\item \textbf{0-10 Scopus/WoS indexed papers}
				\item \textbf{20-35 review papers} (high citation generators)
				\item \textbf{5-10 conference proceedings} (fast publication)
			\end{itemize}
		\end{block}
		\begin{block}{Citation Targets}
			\begin{itemize}
				\item \textbf{1000-2000 total Google Scholar citations}
				\item \textbf{1-10 .gov citations indexing or referencing crossrefs} (Federal Reserve, ERIC, Science.gov, BLS, NIH)
				\item \textbf{10-20 .edu citations} (Harvard, MIT, Stanford, Berkeley)
				\item \textbf{h-index: 15-25+}
			\end{itemize}
		\end{block}
		\begin{block}{Budget}
			\begin{itemize}
				\item \textbf{\$15,000 - \$25,000 total over 2-3 years}
				\item Focus on impact, not expensive journal names
				\item Low-APC journals (\$200-\$800 per paper)
			\end{itemize}
		\end{block}
	\end{frame}
	
	% ============================================================
	% SLIDE 3: BUDGET BREAKDOWN - 25 PAPERS
	% ============================================================
	\begin{frame}{Budget Breakdown - Publishing 25 Papers}
		\scriptsize
		\begin{block}{Cost Per Paper Strategy}
			\centering
			\begin{tabular}{|p{4cm}|p{3cm}|p{3cm}|}
				\hline
				\textbf{Paper Type} & \textbf{Number} & \textbf{Cost Each} \\
				\hline
				Low-APC Journals (\$200-\$400) & 10 & \$2,000-\$4,000 \\
				\hline
				Mid-APC Journals (\$500-\$800) & 8 & \$4,000-\$6,400 \\
				\hline
				No-APC Journals (subscription) & 4 & \$0 \\
				\hline
				Conference Proceedings & 8 & \$300-\$500 each (reg) \\
				\hline
				Book Chapters & 2 & \$0-\$200 \\
				\hline
				Preprints & All & \$0 \\
				\hline
				\textbf{Total} & \textbf{25+} & \textbf{\$10,000-\$20,000} \\
				\hline
			\end{tabular}
		\end{block}
		\begin{block}{Other Costs}
			\begin{itemize}
				\item Professional editing (optional): \$200-\$500 per paper for 5-10 papers = \$5,000-\$10,000
				\item Total with everything: \$10,000 - \$20,000
			\end{itemize}
		\end{block}
	\end{frame}
	
	% ============================================================
	% SLIDE 4: WHY IMPACT > JOURNAL NAME
	% ============================================================
	\begin{frame}{Why Impact Matters More Than Journal Name}
		\scriptsize
		\begin{block}{The Truth About Journal Prestige}
			\begin{itemize}
				\item A paper in a mid-tier journal with 100 citations by gov edu > paper in Nature with 0 citations
				\item USCIS cares about CITATIONS - proof your work is being used
				\item .gov citations from mid-tier journals = as valuable as top-tier
				\item Focus on getting your work READ and CITED, not where it appears
			\end{itemize}
		\end{block}
		\begin{block}{Where to Focus Your Energy - gov angle matters the most}
			\begin{itemize}
				\item \textbf{Content quality:} Solve real problems, address government priorities
				\item \textbf{Distribution:} Upload everywhere (arXiv, SSRN, OSF, ResearchGate)
				\item \textbf{Promotion:} Share on LinkedIn, email to researchers, present at conferences
				\item \textbf{Government alignment:} Topics that federal agencies care about
			\end{itemize}
		\end{block}
	\end{frame}
	
	% ============================================================
	% SLIDE 5: GOVERNMENT ANGLE - THE MOST IMPORTANT FACTOR
	% ============================================================
	\begin{frame}{Government Angle - The Most Important Factor}
		\scriptsize
		\begin{block}{Why Government Impact Matters Most}
			\begin{itemize}
				\item .gov citations are the SINGLE strongest evidence
				\item Federal agencies cite work that serves national interests
				\item One .gov citation can outweigh 50 regular citations
			\end{itemize}
		\end{block}
		\begin{block}{Federal Agencies to Target}
			\centering
			\begin{tabular}{|p{5cm}|p{5cm}|}
				\hline
				\textbf{Agency} & \textbf{Your Research Should Address} \\
				\hline
				Federal Reserve & Financial stability, payment systems, AI in banking \\
				\hline
				Department of Treasury & Crypto regulation, sanctions, fraud detection \\
				\hline
				Department of Labor & Workforce development, job training, AI displacement \\
				\hline
				Department of Education & AI literacy, digital skills, online learning \\
				\hline
				Department of Homeland Security & Cybersecurity, critical infrastructure \\
				\hline
				NIH/PubMed & Health AI, medical data, patient safety \\
				\hline
			\end{tabular}
		\end{block}
	\end{frame}
	
	% ============================================================
	% SLIDE 6: 50 PAPERS - THE VOLUME STRATEGY
	% ============================================================
	\begin{frame}{50 Papers - The High-Volume Strategy}
		\scriptsize
		\begin{block}{Why Volume Matters}
			\begin{itemize}
				\item More papers = more chances to be cited
				\item Each paper can generate 20-50 citations over 2-3 years
				\item 50 papers × 20 average citations = 1000 citations
				\item Volume demonstrates sustained research activity
				\item Manual submission to top indexing like MPRA,EconPaper, SSRN, ERIC, OSF, Sciexplorer
			\end{itemize}
		\end{block}
		\begin{block}{How to Achieve 50 Papers in 2-3 Years}
			\begin{itemize}
				\item \textbf{Write review papers:} 10-15 review papers (each cites 100+ refs, gets many citations)
				\item \textbf{Conference papers:} 15-20 conference proceedings (faster publication)
				\item \textbf{Collaborate:} Each collaboration adds 2-3 co-authored papers per year
				\item \textbf{Extend conference papers:} Conference $\rightarrow$ extended version $\rightarrow$ journal
				\item \textbf{Book chapters:} 5-8 invited chapters (editors reach out after you publish)
			\end{itemize}
		\end{block}
	\end{frame}
	
	% ============================================================
	% SLIDE 7: 1000-2000 CITATIONS - THE STRATEGY
	% ============================================================
	\begin{frame}{1000-2000 Citations - The Strategy}
		\scriptsize
		\begin{block}{Citation Sources to Target}
			\centering
			\begin{tabular}{|p{4cm}|p{5.5cm}|}
				\hline
				\textbf{Source Type} & \textbf{Target Count} \\
				\hline
				Regular journal citations & 800-1500 \\
				\hline
				.gov citations & 1-10 (each high value) \\
				\hline
				.edu citations & 2-20 \\
				\hline
				Conference proceedings citations & 100-200 \\
				\hline
				Preprint citations & 50-100 \\
				\hline
				Self-citations (reasonable) & 200-400 \\
				\hline
			\end{tabular}
		\end{block}
		\begin{block}{Citation Velocity Strategy}
			\begin{itemize}
				\item \textbf{Year 1:} Build foundation (5-10 papers, 0-50 citations)
				\item \textbf{Year 2:} Momentum builds (15-25 papers, 200-500 citations)
				\item \textbf{Year 3:} Exponential growth (30-50 papers, 1000-2000 citations)
				\item Review papers from Year 1 will be heavily cited in Year 2-3
			\end{itemize}
		\end{block}
	\end{frame}
	
	% ============================================================
	% SLIDE 8: REVIEW PAPERS - YOUR CITATION ENGINE
	% ============================================================
	\begin{frame}{Review Papers - Your Citation Engine}
		\scriptsize
		\begin{block}{Why Review Papers Are Critical}
			\begin{itemize}
				\item Review papers get 2-5x MORE citations than original research
				\item Anyone starting research in the field WILL cite your review
				\item Stay relevant and citable for 5-10 years
				\item Can be written without new experiments
			\end{itemize}
		\end{block}
		\begin{block}{How Many Review Papers You Need}
			\begin{itemize}
				\item Target: \textbf{10-15 review papers} over 2-3 years
				\item Each review paper: 150-300 references cited
				\item Each review paper: Expected 100-300 citations over 3 years
				\item 10 review papers × 150 citations = 1500 citations
			\end{itemize}
		\end{block}
		\begin{block}{Review Paper Topics (Government Focus)}
			\begin{itemize}
				\item "AI in Financial Stability: A Systematic Review" (Federal Reserve)
				\item "Cybersecurity for Critical Infrastructure" (DHS)
				\item "Digital Workforce Training: Methods and Outcomes" (DOL)
				\item "AI Literacy in K-12 Education" (Department of Education)
			\end{itemize}
		\end{block}
	\end{frame}
	
	% ============================================================
	% SLIDE 9: RAPID PUBLICATION VENUES
	% ============================================================
	\begin{frame}{Rapid Publication Venues - Build Volume Fast}
		\scriptsize
		\begin{block}{Fastest Publication Options (2-4 months)}
			\begin{itemize}
				\item \textbf{MDPI journals:} 4-6 weeks from submission to publication
				\item \textbf{Elsevier open access:} 6-12 weeks
				\item \textbf{IEEE conferences:} 3-4 months from submission to presentation
				\item \textbf{Springer conference proceedings:} 3-5 months
			\end{itemize}
		\end{block}
		\begin{block}{Low-APC Rapid Journals (Under \$500)}
			\begin{itemize}
				\item MDPI: Some journals have low APCs (\$400-\$800) or waivers
				\item Hindawi: \$800-\$1,200 (negotiate waivers)
				\item PeerJ: \$399-\$1,095 (tiered pricing)
				\item Directory of Open Access Journals: Filter by APC \$0-\$500
			\end{itemize}
		\end{block}
	\end{frame}
	
	% ============================================================
	% SLIDE 10: ZERO-APC PUBLICATION STRATEGY
	% ============================================================
	\begin{frame}{Zero-APC Publication Strategy}
		\scriptsize
		\begin{block}{No-Cost Publication Options}
			\begin{itemize}
				\item \textbf{Traditional subscription journals:} Author pays nothing (e.g., IEEE Transactions, many Elsevier subscription journals)
				\item \textbf{University journals:} Many university presses charge no APC
				\item \textbf{Society journals:} IEEE, ACM, INFORMS member rates often \$0
				\item \textbf{Conference proceedings:} Registration fee only (\$300-\$800)
				\item \textbf{Book chapters:} Usually no fee (invited)
				\item \textbf{Preprints:} Always free (arXiv, SSRN, OSF)
			\end{itemize}
		\end{block}
		\begin{block}{How to Find No-APC Journals}
			\begin{enumerate}
				\item Go to Scopus Sources
				\item Search your field
				\item Check "Open Access" filter - select "All Open Access" to see options
				\item Or select "Hybrid" journals (subscription with OA option - you can choose subscription)
				\item Email editor to confirm no fee for subscription track
			\end{enumerate}
		\end{block}
	\end{frame}
	
	% ============================================================
	% SLIDE 11: CONFERENCE PROCEEDINGS - VOLUME BUILDER
	% ============================================================
	\begin{frame}{Conference Proceedings - High-Volume Strategy}
		\scriptsize
		\begin{block}{Why Conferences for Volume}
			\begin{itemize}
				\item Faster than journals (3-5 months vs 12+ months)
				\item Lower cost (registration \$300-\$800 vs APC \$1,000+)
				\item Can publish 3-5 conference papers per year
				\item Easy to extend into journal papers later
			\end{itemize}
		\end{block}
		\begin{block}{Conference Proceedings Indexed in Scopus}
			\begin{itemize}
				\item \textbf{IEEE Xplore:} All IEEE conferences (strongest)
				\item \textbf{Springer LNCS/LNAI/CCIS:} High-quality proceedings
				\item \textbf{ACM Digital Library:} ACM conferences
				\item \textbf{Elsevier Procedia:} Procedia Computer Science, Procedia Economics
				\item \textbf{CEUR Workshop Proceedings:} Free, Scopus indexed for some
			\end{itemize}
		\end{block}
		\begin{block}{Conference Strategy}
			\begin{itemize}
				\item Submit 2-3 conference papers per year
				\item Each conference = 1 publication, 5-20 eventual citations
				\item Present = networking for collaborations
			\end{itemize}
		\end{block}
	\end{frame}
	
	% ============================================================
	% SLIDE 12: COLLABORATION FOR VOLUME
	% ============================================================
	\begin{frame}{Collaboration - Force Multiplier for Volume}
		\scriptsize
		\begin{block}{Why Collaborate}
			\begin{itemize}
				\item Each collaboration adds 2-5 co-authored papers per year
				\item You don't have to write everything yourself
				\item Shared workload = more publications for everyone
			\end{itemize}
		\end{block}
		\begin{block}{How to Find Collaborators}
			\begin{itemize}
				\item \textbf{ResearchGate:} Message researchers in your field
				\item \textbf{LinkedIn:} Connect with academics and industry researchers
				\item \textbf{Conferences:} Meet potential collaborators in person
				\item \textbf{Twitter/X:} Follow and engage with researchers
				\item \textbf{Former colleagues:} Reach out to past coworkers
			\end{itemize}
		\end{block}
		\begin{block}{Collaboration Structure}
			\begin{itemize}
				\item You lead: Write first draft, handle submission (you get first author)
				\item Partner leads: They write, you contribute (you get co-author)
				\item Multiple small collaborations = many co-authored papers
			\end{itemize}
		\end{block}
	\end{frame}
	
	% ============================================================
	% SLIDE 13: PREPRINTS - CITATIONS BEFORE PUBLICATION
	% ============================================================
	\begin{frame}{Preprints - Citations Before Journal Publication}
		\scriptsize
		\begin{block}{Why Preprints Are Essential}
			\begin{itemize}
				\item Citations start the DAY you upload (not 12 months later)
				\item Google Scholar indexes preprints
				\item Preprints get DOI and become citable immediately
				\item Can be cited by government reports before journal publication
			\end{itemize}
		\end{block}
		\begin{block}{Where to Upload Every Paper}
			\begin{enumerate}
				\item \textbf{arXiv.org} (CS, physics, math, finance)
				\item \textbf{SSRN} (economics, finance, law)
				\item \textbf{OSF Preprints} (interdisciplinary, DOI from OSF)
				\item \textbf{ResearchGate} (upload full text)
			\end{enumerate}
		\end{block}
		\begin{block}{Preprint Citation Timeline Example}
			\begin{itemize}
				\item Day 0: Upload to arXiv
				\item Day 30: Paper indexed in Google Scholar
				\item Day 60: First citation (someone read your preprint)
				\item Month 6: Journal publication (citations already started)
				\item Month 12: 20-50 citations (many from preprint version)
			\end{itemize}
		\end{block}
	\end{frame}
	
	% ============================================================
	% SLIDE 14: .GOV CITATIONS - FEDERAL RESERVE
	% ============================================================
	\begin{frame}{.gov Citations - Federal Reserve (FEDS Series)}
		\scriptsize
		\begin{block}{Targeting Federal Reserve Citations}
			\begin{itemize}
				\item Federal Reserve economists read SSRN and arXiv preprints
				\item Your preprint can be cited BEFORE journal publication
				\item FEDS papers have DOIs and are .gov domain
			\end{itemize}
		\end{block}
		\begin{block}{Topics That Get Fed Citations}
			\begin{itemize}
				\item Machine learning for financial stability monitoring
				\item Payment systems and FedNow adoption
				\item Climate risk and financial institutions
				\item AI in banking supervision and compliance
				\item Cryptocurrency and digital asset regulation
				\item Stress testing methodologies
			\end{itemize}
		\end{block}
		\begin{block}{How to Maximize Fed Citation Chance}
			\begin{enumerate}
				\item Upload to SSRN (Fed economists use SSRN heavily)
				\item Include "Federal Reserve" or "financial stability" in keywords
				\item Cite Fed papers in your work (they notice)
				\item Email your SSRN link to Fed economists (appropriate ones)
			\end{enumerate}
		\end{block}
	\end{frame}
	
	% ============================================================
	% SLIDE 15: .GOV CITATIONS - ERIC (DEPARTMENT OF EDUCATION)
	% ============================================================
	\begin{frame}{.gov Citations - ERIC (Department of Education)}
		\scriptsize
		\begin{block}{What is ERIC?}
			\begin{itemize}
				\item Education Resources Information Center
				\item .gov domain of U.S. Department of Education
				\item Indexes education, workforce, training research
			\end{itemize}
		\end{block}
		\begin{block}{How Papers Get into ERIC}
			\begin{enumerate}
				\item Publish in ERIC-indexed journal
				\item Journal automatically sends content to ERIC
				\item Your paper appears on ERIC with .gov URL
				\item Now your work is discoverable through .gov portal
			\end{enumerate}
		\end{block}
		\begin{block}{Topics for ERIC Indexing}
			\begin{itemize}
				\item AI literacy and digital skills training
				\item Workforce development and retraining programs
				\item Online learning effectiveness research
				\item Veteran education and job training
				\item STEM education and diversity
				\item Community college workforce programs
			\end{itemize}
		\end{block}
	\end{frame}
	
	% ============================================================
	% SLIDE 16: .GOV CITATIONS - SCIENCE.GOV
	% ============================================================
	\begin{frame}{.gov Citations - Science.gov}
		\scriptsize
		\begin{block}{What is Science.gov?}
			\begin{itemize}
				\item Official U.S. government search portal for science
				\item Searches 60+ federal databases
				\item Managed by Department of Energy
			\end{itemize}
		\end{block}
		\begin{block}{Pathways to Science.gov}
			\begin{itemize}
				\item \textbf{PubMed} (NIH) $\rightarrow$ Science.gov
				\item \textbf{ERIC} (Dept of Education) $\rightarrow$ Science.gov
				\item \textbf{NASA ADS} $\rightarrow$ Science.gov
				\item \textbf{DOE OSTI} $\rightarrow$ Science.gov
				\item \textbf{USDA} $\rightarrow$ Science.gov
			\end{itemize}
		\end{block}
		\begin{block}{Strategy for Science.gov Indexing}
			\begin{itemize}
				\item Publish in PubMed-indexed journals for health/AI topics
				\item Publish in ERIC-indexed journals for education/workforce
				\item Upload to DOE OSTI if your research has energy relevance
				\item Once indexed, your work is searchable on .gov domain
			\end{itemize}
		\end{block}
	\end{frame}
	
	% ============================================================
	% SLIDE 17: .GOV CITATIONS - BUREAU OF LABOR STATISTICS
	% ============================================================
	\begin{frame}{.gov Citations - Bureau of Labor Statistics}
		\scriptsize
		\begin{block}{Targeting BLS Citations}
			\begin{itemize}
				\item BLS publishes Monthly Labor Review (MLR)
				\item Also publishes working papers and reports
				\item Citations from BLS = .gov domain evidence
			\end{itemize}
		\end{block}
		\begin{block}{Topics for BLS Citations}
			\begin{itemize}
				\item Employment projections and labor market trends
				\item Impact of AI on employment and wages
				\item Workforce demographics and participation rates
				\item Occupational outlook and skills demand
				\item Gig economy and alternative work arrangements
			\end{itemize}
		\end{block}
		\begin{block}{How to Get Cited by BLS}
			\begin{enumerate}
				\item Publish on labor economics topics
				\item Upload to SSRN (BLS researchers use SSRN)
				\item Submit comments to BLS RFIs on Regulations.gov
				\item Cite BLS data in your papers (they notice)
			\end{enumerate}
		\end{block}
	\end{frame}
	
	% ============================================================
	% SLIDE 18: .GOV CITATIONS - NIH/PUBMED
	% ============================================================
	\begin{frame}{.gov Citations - NIH / PubMed}
		\scriptsize
		\begin{block}{Targeting NIH Citations}
			\begin{itemize}
				\item PubMed is .gov domain (nih.gov)
				\item Citations from NIH-funded research are .gov
				\item Health-related AI research is high priority
			\end{itemize}
		\end{block}
		\begin{block}{Topics for PubMed Citations}
			\begin{itemize}
				\item AI for clinical decision support
				\item Machine learning for patient safety
				\item Digital health and telemedicine
				\item Health data privacy and security
				\item Medical image analysis
				\item Public health surveillance
			\end{itemize}
		\end{block}
		\begin{block}{How to Get PubMed Citations}
			\begin{enumerate}
				\item Publish in PubMed-indexed journals
				\item Upload to medRxiv preprint server
				\item Collaborate with NIH-funded researchers
				\item Focus on health applications of your methods
			\end{enumerate}
		\end{block}
	\end{frame}
	
	% ============================================================
	% SLIDE 19: .EDU CITATIONS - TOP UNIVERSITIES
	% ============================================================
	\begin{frame}{.edu Citations - Top US Universities}
		\scriptsize
		\begin{block}{Why .edu Citations Are Valuable}
			\begin{itemize}
				\item Verifiable .edu domain
				\item Demonstrates academic recognition in the U.S.
				\item Harvard, MIT, Stanford citations are strongest
			\end{itemize}
		\end{block}
		\begin{block}{How to Get .edu Citations}
			\begin{enumerate}
				\item Publish in reputable journals (Scopus/WoS)
				\item Upload to institutional repositories (Harvard DASH, MIT DSpace)
				\item Present at university seminars (virtual or in-person)
				\item Collaborate with US university researchers
				\item Create course materials that professors adopt
			\end{enumerate}
		\end{block}
		\begin{block}{Tracking .edu Citations}
			\begin{itemize}
				\item Search Google Scholar with this query every 3 months
				\item Save screenshots of results with .edu URLs
				\item Document citations from specific universities
			\end{itemize}
		\end{block}
	\end{frame}
	
	
	
	% ============================================================
	% SLIDE 19.5: EBSCO HOST - MAJOR COMMERCIAL DATABASE
	% ============================================================
	\begin{frame}{EBSCO Host - Massive Academic Reach}
		\scriptsize
		\begin{block}{What is EBSCO?}
			\begin{itemize}
				\item Largest academic database aggregator worldwide
				\item Serves 10,000+ universities globally
				\item Indexes 10,000+ journals across all disciplines
				\item Content appears in university library search results
			\end{itemize}
		\end{block}
		\begin{block}{Why EBSCO Citations Matter}
			\begin{itemize}
				\item Your paper appears when students \& researchers search their university library
				\item More discoverability = more citations
				\item Many Scopus-indexed journals are ALSO in EBSCO
				\item EBSCO discovery services power university search portals
			\end{itemize}
		\end{block}
		\begin{block}{How to Get Indexed in EBSCO}
			\begin{itemize}
				\item Publish in journals that EBSCO indexes (check Journal Finder)
				\item Many low-APC journals are EBSCO-indexed
				\item EBSCO indexes conference proceedings and book chapters too
				\item Ask your journal editor: "Is this journal in EBSCO?"
			\end{itemize}
		\end{block}
	\end{frame}
	
	
	% ============================================================
	% SLIDE 19.6: PROQUEST - DISSERTATIONS AND BEYOND
	% ============================================================
	\begin{frame}{ProQuest - Beyond Traditional Journals}
		\scriptsize
		\begin{block}{What is ProQuest?}
			\begin{itemize}
				\item Major academic database alongside EBSCO
				\item Indexes dissertations, journals, newspapers, conference papers
				\item Available in 99\% of US university libraries
				\item ProQuest Central is most widely used
			\end{itemize}
		\end{block}
		\begin{block}{ProQuest Citation Benefits}
			\begin{itemize}
				\item Citations from .edu domains via university library access
				\item Dissertations citing your work = .edu citations
				\item Conference proceedings indexed in ProQuest get more visibility
				\item Government reports indexed here too
			\end{itemize}
		\end{block}
		\begin{block}{How to Get into ProQuest}
			\begin{itemize}
				\item Conference proceedings often automatically go to ProQuest
				\item Dissertations that cite you are in PQDT (ProQuest Dissertations)
				\item Publish in ProQuest-indexed journals
				\item Upload to institutional repositories (university libraries submit)
			\end{itemize}
		\end{block}
	\end{frame}
	
	% ============================================================
	% SLIDE 20: GETTING STARTED - FIRST 30 DAYS
	% ============================================================
	\begin{frame}{First 30 Days - Setting Up}
		\scriptsize
		\begin{block}{Week 1: Infrastructure}
			\begin{itemize}
				\item Create Google Scholar profile (public)
				\item Create ORCID profile
				\item Create ResearchGate account
				\item Install Zotero
				\item Set up Overleaf or TeXstudio
				\item Create GitHub account
			\end{itemize}
		\end{block}
		\begin{block}{Week 2: Topic Selection}
			\begin{itemize}
				\item Identify 5 research topics aligned with government priorities
				\item Search Regulations.gov for RFIs in your area
				\item Read 10-20 papers on each topic
				\item Choose 3 topics for your first papers
			\end{itemize}
		\end{block}
		\begin{block}{Week 3-4: First Paper}
			\begin{itemize}
				\item Write first review paper (2000-4000 words)
				\item Cite 50-100 references in Zotero
				\item Share draft with 1-2 collaborators
				\item Submit to mid-tier journal
			\end{itemize}
		\end{block}
	\end{frame}
	
	% ============================================================
	% SLIDE 21: MONTHS 1-6 - FIRST 5 PAPERS
	% ============================================================
	\begin{frame}{Months 1-6 - First 5 Papers}
		\scriptsize
		\begin{block}{Month 1-2: Paper 1 (Review Paper)}
			\begin{itemize}
				\item Write and submit first review paper
				\item Target: Mid-tier journal, APC \$500-\$800 or no-APC
				\item Upload to arXiv/SSRN immediately
			\end{itemize}
		\end{block}
		\begin{block}{Month 2-3: Paper 2 (Conference Paper)}
			\begin{itemize}
				\item Find conference with proceedings in Springer/IEEE
				\item Write 6-8 page conference paper
				\item Submit to conference (deadline typically 3-4 months before)
			\end{itemize}
		\end{block}
		\begin{block}{Month 3-4: Paper 3 (Short Original Research)}
			\begin{itemize}
				\item Extend Paper 2 into journal paper (add 30\% new content)
				\item Submit to low-APC journal (\$200-\$400)
			\end{itemize}
		\end{block}
		\begin{block}{Month 4-5: Paper 4 (Second Review Paper)}
			\begin{itemize}
				\item Different topic from Paper 1
				\item Submit to different journal
			\end{itemize}
		\end{block}
		\begin{block}{Month 5-6: Paper 5 (Collaboration Paper)}
			\begin{itemize}
				\item Reach out to 1-2 collaborators on ResearchGate
				\item Co-author a paper on shared interest
			\end{itemize}
		\end{block}
	\end{frame}
	
	% ============================================================
	% SLIDE 22: MONTHS 7-12 - PAPERS 6-15
	% ============================================================
	\begin{frame}{Months 7-12 - Papers 6-15}
		\scriptsize
		\begin{block}{Accelerating Publication Rate}
			\begin{itemize}
				\item \textbf{By month 6:} You have 5 papers published or under review
				\item \textbf{Month 7-9:} Submit 3-4 more papers (mix of reviews and originals)
				\item \textbf{Month 10-12:} Submit 3-4 more papers
				\item \textbf{Total by month 12:} 12-15 papers
			\end{itemize}
		\end{block}
		\begin{block}{Peer Review Begins}
			\begin{itemize}
				\item By month 6-8: First peer review invitations arrive
				\item Complete 2-3 reviews by month 12
				\item Add to Web of Science reviewer profile
			\end{itemize}
		\end{block}
		\begin{block}{First Citations Appear}
			\begin{itemize}
				\item By month 9-12: First citations on Google Scholar
				\item Preprints from Month 1-2 are now being cited
				\item Target: 20-50 citations by month 12
			\end{itemize}
		\end{block}
	\end{frame}
	
	% ============================================================
	% SLIDE 23: YEAR 2 - PAPERS 16-35
	% ============================================================
	\begin{frame}{Year 2 - Papers 16-35}
		\scriptsize
		\begin{block}{Publication Cadence}
			\begin{itemize}
				\item \textbf{Month 13-18:} Submit 8-10 papers
				\item \textbf{Month 19-24:} Submit 8-10 papers
				\item \textbf{Total Year 2:} 15-20 new papers
				\item \textbf{Cumulative:} 25-35 total papers
			\end{itemize}
		\end{block}
		\begin{block}{Citation Growth Accelerates}
			\begin{itemize}
				\item Papers from Year 1 now have 20-50 citations each
				\item Review papers from Year 1 are heavily cited (50-150 citations)
				\item Total citations: 300-800 by end of Year 2
			\end{itemize}
		\end{block}
		\begin{block}{Peer Review Portfolio Grows}
			\begin{itemize}
				\item Complete 15-25 reviews in Year 2
				\item Review for 5-8 distinct journals
				\item Web of Science reviewer recognition established
			\end{itemize}
		\end{block}
	\end{frame}
	
	% ============================================================
	% SLIDE 24: YEAR 3 - PAPERS 36-50
	% ============================================================
	\begin{frame}{Year 3 - Papers 36-50}
		\scriptsize
		\begin{block}{Final Push to 50 Papers}
			\begin{itemize}
				\item \textbf{Month 25-30:} Submit 8-10 papers
				\item \textbf{Month 31-36:} Submit 8-10 papers
				\item \textbf{Total Year 3:} 15-20 new papers
				\item \textbf{Cumulative:} 45-55 total papers
			\end{itemize}
		\end{block}
		\begin{block}{Citation Target Achieved}
			\begin{itemize}
				\item Year 1 papers: 100-300 citations each (review papers)
				\item Year 2 papers: 30-100 citations each
				\item Year 3 papers: 10-50 citations each
				\item \textbf{Total: 1000-2000+ citations}
			\end{itemize}
		\end{block}
		\begin{block}{Government Citations}
			\begin{itemize}
				\item By Year 3: 5-15 .gov citations (Fed, ERIC, BLS, etc.)
				\item By Year 3: 15-30 .edu citations
				\item Portfolio complete and ready for filing
			\end{itemize}
		\end{block}
	\end{frame}
	
	% ============================================================
	% SLIDE 25: SAMPLE PAPER TOPIC - FINANCIAL STABILITY
	% ============================================================
	\begin{frame}{Sample Paper Topics - Financial Stability}
		\scriptsize
		\begin{block}{Review Papers (10-15 of these)}
			\begin{enumerate}
				\item "Machine Learning for Financial Stability: A Systematic Review"
				\item "AI in Banking Supervision: Methods and Applications"
				\item "Cryptocurrency Regulation: A Review of Global Approaches"
				\item "Stress Testing Methodologies: Evolution and Future Directions"
				\item "Payment Systems Security: Threats and Countermeasures"
			\end{enumerate}
		\end{block}
		\begin{block}{Original Research Papers (15-25 of these)}
			\begin{enumerate}
				\item "Detecting Anomalies in Real-Time Payment Systems Using Graph Neural Networks"
				\item "Forecasting Bank Runs with Social Media Sentiment Analysis"
				\item "Explainable AI for Anti-Money Laundering Compliance"
				\item "Blockchain for Cross-Border Payments: Efficiency Analysis"
				\item "Machine Learning for Credit Risk Assessment in Community Banks"
			\end{enumerate}
		\end{block}
	\end{frame}
	
	% ============================================================
	% SLIDE 26: SAMPLE PAPER TOPIC - WORKFORCE DEVELOPMENT
	% ============================================================
	\begin{frame}{Sample Paper Topics - Workforce Development}
		\scriptsize
		\begin{block}{Review Papers}
			\begin{enumerate}
				\item "AI and Workforce Displacement: A Systematic Review of Retraining Programs"
				\item "Digital Literacy Frameworks for Adult Learners"
				\item "Online Learning Effectiveness for Workforce Training"
				\item "Veteran Transition to Tech Careers: Best Practices"
				\item "Community College Workforce Programs: Models and Outcomes"
			\end{enumerate}
		\end{block}
		\begin{block}{Original Research Papers}
			\begin{enumerate}
				\item "Evaluating AI Literacy Curriculum for Community College Students"
				\item "Predicting Workforce Training Outcomes Using Machine Learning"
				\item "Gamification for Digital Skills Training: Randomized Controlled Trial"
				\item "Employer Perspectives on AI Credentials and Certifications"
				\item "Economic Impact of Digital Upskilling Programs on Local Communities"
			\end{enumerate}
		\end{block}
	\end{frame}
	
	% ============================================================
	% SLIDE 27: SAMPLE PAPER TOPIC - CYBERSECURITY
	% ============================================================
	\begin{frame}{Sample Paper Topics - Cybersecurity}
		\scriptsize
		\begin{block}{Review Papers}
			\begin{enumerate}
				\item "AI for Cybersecurity: A Systematic Review of Detection Methods"
				\item "Critical Infrastructure Protection: Frameworks and Gaps"
				\item "Zero-Trust Architecture: Implementation Challenges and Solutions"
				\item "Ransomware Trends and Defense Strategies: 2020-2025"
				\item "Supply Chain Cybersecurity: Risk Assessment Methodologies"
			\end{enumerate}
		\end{block}
		\begin{block}{Original Research Papers}
			\begin{enumerate}
				\item "Detecting Phishing Attacks Using Large Language Models"
				\item "Anomaly Detection in Industrial Control Systems with Autoencoders"
				\item "Federated Learning for Collaborative Threat Intelligence"
				\item "Automated Vulnerability Assessment Using Graph Neural Networks"
				\item "Privacy-Preserving Machine Learning for Healthcare Data"
			\end{enumerate}
		\end{block}
	\end{frame}
	
	% ============================================================
	% SLIDE 28: SAMPLE PAPER TOPIC - AI IN GOVERNMENT
	% ============================================================
	\begin{frame}{Sample Paper Topics - AI in Government}
		\scriptsize
		\begin{block}{Review Papers}
			\begin{enumerate}
				\item "AI Governance Frameworks: A Comparative Analysis of Federal Agencies"
				\item "Algorithmic Bias Assessment Methods for Government Systems"
				\item "Explainable AI for Regulatory Compliance"
				\item "AI Procurement Best Practices for Public Sector"
				\item "Privacy-Preserving AI for Government Data Sharing"
			\end{enumerate}
		\end{block}
		\begin{block}{Original Research Papers}
			\begin{enumerate}
				\item "Automating FOIA Redaction Using NLP"
				\item "AI for Regulatory Comment Analysis at Scale"
				\item "Predictive Models for Government Program Outcomes"
				\item "Chatbots for Citizen Service Delivery: Effectiveness Study"
				\item "AI-Assisted Policy Impact Assessment: Case Studies"
			\end{enumerate}
		\end{block}
	\end{frame}
	
	% ============================================================
	% SLIDE 29: MAXIMIZING .GOV CITATIONS - ACTION PLAN
	% ============================================================
	\begin{frame}{Maximizing .gov Citations - Action Plan}
		\scriptsize
		\begin{block}{Monthly Actions for .gov Citations}
			\begin{itemize}
				\item \textbf{Each month:} Check Regulations.gov for new RFIs in your area
				\item \textbf{Each month:} Submit 1 comment to an RFI citing your own research
				\item \textbf{Each quarter:} Search for new .gov citations to your papers
				\item \textbf{Each quarter:} Email your SSRN/arXiv links to agency researchers
			\end{itemize}
		\end{block}
		\begin{block}{Agency Outreach Strategy}
			\begin{enumerate}
				\item Identify 10-20 researchers at Federal Reserve, BLS, ERIC, NIH
				\item Find their email (agency websites, Google Scholar)
				\item Send professional email: "I noticed your work on [topic]. I have a preprint you might find relevant: [link]"
				\item No requests, just sharing - they may cite if they find it useful
			\end{enumerate}
		\end{block}
	\end{frame}
	
	% ============================================================
	% SLIDE 30: CITATION TRACKING DASHBOARD
	% ============================================================
% ============================================================
% SLIDE 30: CITATION TRACKING DASHBOARD
% ============================================================
\begin{frame}{Citation Tracking Dashboard}
	\scriptsize
	\begin{block}{What to Track Monthly}
		\centering
		\begin{tabular}{|p{4cm}|p{5.5cm}|}
			\hline
			\textbf{Metric} & \textbf{How to Track} \\
			\hline
			Total Google Scholar citations & Google Scholar profile \\
			\hline
			Monthly new citations & Google Scholar "Cited by" alerts \\
			\hline
			.gov citations & Google Scholar search: site:.gov "your name" \\
			\hline
			.edu citations & Google Scholar search: site:.edu "your name" \\
			\hline
			h-index & Google Scholar profile \\
			\hline
			Peer reviews completed & Web of Science / ORCID \\
			\hline
			Preprint downloads & arXiv/SSRN/OSF stats \\
			\hline
		\end{tabular}
	\end{block}
	\begin{block}{Sample Tracking Sheet}
		\centering
		\begin{tabular}{|c|c|c|c|c|c|}
			\hline
			\textbf{Month} & \textbf{Papers} & \textbf{Citations} & \textbf{.gov} & \textbf{.edu} & \textbf{Reviews} \\
			\hline
			Month 6 & 8 & 25 & 0 & 2 & 3 \\
			\hline
			Month 12 & 15 & 80 & 1 & 5 & 8 \\
			\hline
			Month 18 & 25 & 300 & 3 & 12 & 15 \\
			\hline
			Month 24 & 35 & 800 & 5 & 25 & 25 \\
			\hline
			Month 30 & 45 & 1500 & 7 & 40 & 35 \\
			\hline
			Month 36 & 50 & 2000 & 8 & 50 & 50 \\
			\hline
		\end{tabular}
	\end{block}
\end{frame}
	
	% ============================================================
	% SLIDE 31: LOW-APC JOURNAL LIST
	% ============================================================
	\begin{frame}{Low-APC Journal List (\$200-\$800)}
		\scriptsize
		\begin{block}{Verified Low-APC Journals (Check Current Status)}
			\begin{itemize}
				\item \textbf{Data in Brief} (Elsevier) - APC \$500
				\item \textbf{Software Impacts} (Elsevier) - APC \$600
				\item \textbf{Array} (Elsevier) - APC \$700
				\item \textbf{SN Computer Science} (Springer) - APC \$800 (waivers available)
				\item \textbf{Journal of Engineering and Applied Science} (Springer) - APC \$600
				\item \textbf{PeerJ Computer Science} - APC \$399-\$1,095 (tiered)
				\item \textbf{Heliyon} (Elsevier) - APC \$1,250 (higher but negotiate)
			\end{itemize}
		\end{block}
		
	\end{frame}
	
	% ============================================================
	% SLIDE 32: AVOID EXPENSIVE JOURNALS
	% ============================================================
	\begin{frame}{Avoid Expensive Journals - Low ROI for Citations}
		\scriptsize
		\begin{block}{Journals to Avoid (High APC, Low Citation ROI)}
			\begin{itemize}
				\item \textcolor{red}{Nature Communications:} APC \$6,000+ (not worth for citation volume)
				\item \textcolor{red}{Scientific Reports:} APC \$2,500+ (too expensive for budget)
				\item \textcolor{red}{Frontiers journals:} APC \$2,000-\$3,500 (expensive)
				\item \textcolor{red}{PLoS ONE:} APC \$1,800 (borderline - only if waiver obtained)
				\item \textcolor{red}{MDPI top journals:} APC \$1,800-\$2,500 (too expensive for volume)
			\end{itemize}
		\end{block}
		\begin{block}{Why to Avoid}
			\begin{itemize}
				\item Your budget is \$15-25K for 50 papers = average \$300-\$500 per paper
				\item \$2,000 per paper × 50 = \$100,000 (impossible)
				\item Focus on low-APC and no-APC journals
				\item Citations come from CONTENT, not journal prestige
			\end{itemize}
		\end{block}
	\end{frame}
	
	% ============================================================
	% SLIDE 33: PEER REVIEW - 50+ REVIEWS STRATEGY
	% ============================================================
	\begin{frame}{Peer Review - 50+ Reviews Strategy}
		\scriptsize
		\begin{block}{Target: 50+ Peer Reviews over 3 Years}
			\begin{itemize}
				\item Year 1: 5-10 reviews (starting out)
				\item Year 2: 15-25 reviews (established reviewer)
				\item Year 3: 25-35 reviews (senior reviewer)
			\end{itemize}
		\end{block}
		\begin{block}{How to Get Invited}
			\begin{enumerate}
				\item Publish 3-5 papers in a journal (editors notice)
				\item Register as reviewer on ScholarOne and Editorial Manager
				\item Email journal editors: "I have published in your journal and am available to review"
				\item Accept 50-70\% of invitations (decline if overloaded)
			\end{enumerate}
		\end{block}
		\begin{block}{Review Documentation}
			\begin{itemize}
				\item Save all invitation emails
				\item Save completion confirmations
				\item Web of Science reviewer profile (automatic tracking)
				\item ORCID reviewer recognition
			\end{itemize}
		\end{block}
	\end{frame}
	
	% ============================================================
	% SLIDE 34: AWARDS - BUILDING RECOGNITION
	% ============================================================
	\begin{frame}{Awards - Building Recognition on a Budget}
		\scriptsize
		\begin{block}{Free/Low-Cost Award Opportunities}
			\begin{itemize}
				\item \textbf{Best Paper Awards} at conferences you attend (included in registration)
				\item \textbf{Young Researcher Awards} from professional societies (often free to apply)
				\item \textbf{Travel Grants} (award for attending - still counts as award)
				\item \textbf{Outstanding Reviewer Awards} from journals (free, based on your reviews)
				\item \textbf{Teaching Awards} from your employer
			\end{itemize}
		\end{block}
		\begin{block}{Awards to Target}
			\begin{itemize}
				\item IEEE Outstanding Young Professional (nomination-based)
				\item ACM Student Research Competition (if affiliated with university)
				\item Conference Best Paper Award (submit your best work)
				\item Journal Outstanding Reviewer Award (complete 10+ reviews for one journal)
			\end{itemize}
		\end{block}
	\end{frame}
	
	% ============================================================
	% SLIDE 35: BOOK CHAPTERS - ADDING 5-10 CHAPTERS
	% ============================================================
	\begin{frame}{Book Chapters - Adding 5-10 Chapters}
		\scriptsize
		\begin{block}{How to Get Book Chapter Invitations}
			\begin{itemize}
				\item Publish 5-10 journal papers first
				\item Editors of edited volumes will find you on Google Scholar
				\item They email: "Would you like to contribute a chapter on topic X"
				\item Accept 50\% of invitations
			\end{itemize}
		\end{block}
		\begin{block}{Cost for Book Chapters}
			\begin{itemize}
				\item Most book chapters have NO APC
				\item Some publishers charge \$200-\$500 (negotiate or decline)
				\item Target: 5-10 chapters over 3 years at \$0 cost
			\end{itemize}
		\end{block}
		\begin{block}{Publishers with Free Book Chapters}
			\begin{itemize}
				\item Springer (usually no fee for invited chapters)
				\item Elsevier (check for fee)
				\item Wiley (check for fee)
				\item IGI Global (may have fee - avoid unless invited with waiver)
			\end{itemize}
		\end{block}
	\end{frame}
	
	% ============================================================
	% SLIDE 36: PREPRINT SERVERS - EVERY PAPER GOES HERE
	% ============================================================
	\begin{frame}{Preprint Servers - Every Paper Must Go Here}
		\scriptsize
		\begin{block}{Upload Every Paper to ALL These Servers}
			\begin{enumerate}
				\item \textbf{arXiv.org} (for CS, physics, math, finance, statistics)
				\item \textbf{SSRN} (for economics, finance, law, social sciences)
				\item \textbf{OSF Preprints} (for interdisciplinary, gets DOI)
				\item \textbf{ResearchGate} (full text upload)
			\end{enumerate}
		\end{block}
		\begin{block}{Why This Matters for Citations}
			\begin{itemize}
				\item Each server = additional discovery pathway
				\item Each server = potential new reader = potential citation
				\item 4 servers × 50 papers = 200 distribution points
				\item Citations start the day you upload, not 12 months later
			\end{itemize}
		\end{block}
		\begin{block}{Preprint Citation Evidence for Portfolio}
			\begin{itemize}
				\item Screenshot of arXiv/SSRN download counts
				\item Citation of your preprint in other papers
				\item Google Scholar shows preprint citations
			\end{itemize}
		\end{block}
	\end{frame}
	

	
	% ============================================================
	% SLIDE 38: TWITTER/X FOR RESEARCH NETWORKING
	% ============================================================
	\begin{frame}{Twitter/X - Research Networking}
		\scriptsize
		\begin{block}{Why Twitter/X for Researchers}
			\begin{itemize}
				\item Active academic community (hashtags: \#AcademicTwitter, \#AI, \#MachineLearning)
				\item Journal editors and conference organizers are on Twitter
				\item Preprints shared here get more citations
			\end{itemize}
		\end{block}
		\begin{block}{Twitter Strategy}
			\begin{itemize}
				\item Share every paper with 1-2 sentence summary
				\item Use 3-5 relevant hashtags
				\item Tag journal/conference account if applicable
				\item Retweet and engage with other researchers' papers
			\end{itemize}
		\end{block}
	\end{frame}
	
	% ============================================================
	% SLIDE 39: ZOTERO FOR LITERATURE REVIEW - SCALE TO 1000+ REFERENCES
	% ============================================================
	\begin{frame}{Zotero - Literature Review at Scale}
		\scriptsize
		\begin{block}{Why Zotero for 50 Papers}
			\begin{itemize}
				\item Each review paper cites 100-300 references
				\item 10-15 review papers = 1500-4500 references managed
				\item Zotero organizes everything automatically
			\end{itemize}
		\end{block}
		\begin{block}{Zotero Setup for Volume}
			\begin{itemize}
				\item Create collections for each paper
				\item Use tags: "\#review", "\#methodology", "\#citation"
				\item Install Better BibTeX for LaTeX export
				\item Use browser connector to save papers instantly
			\end{itemize}
		\end{block}
		\begin{block}{Literature Review Workflow}
			\begin{enumerate}
				\item Search Google Scholar for 100-200 relevant papers
				\item Save to Zotero with one click
				\item Export .bib file for your LaTeX paper
				\item Write review paper citing these references
			\end{enumerate}
		\end{block}
	\end{frame}
	
	% ============================================================
	% SLIDE 40: LATEX FOR 50 PAPERS - TEMPLATE SYSTEM
	% ============================================================
% ============================================================
% SLIDE 40: LATEX FOR 50 PAPERS - TEMPLATE SYSTEM
% ============================================================
\begin{frame}{LaTeX - Template System for 50 Papers}
	\scriptsize
	\begin{block}{Why LaTeX for Volume}
		\begin{itemize}
			\item Create templates - reuse for every paper
			\item Automatic bibliography (BibTeX) - no manual formatting
			\item Consistent formatting across 50 papers
			\item Free (Overleaf or TeXstudio)
		\end{itemize}
	\end{block}
	\begin{block}{Template Structure}
		\centering
		\begin{tabular}{ll}
			\multicolumn{2}{l}{\texttt{paper\_template.tex}} \\
			\texttt{$\vdash$ preamble.tex} & (packages, commands) \\
			\texttt{$\vdash$ title.tex} & (title, authors, abstract) \\
			\texttt{$\vdash$ introduction.tex} & \\
			\texttt{$\vdash$ literature\_review.tex} & \\
			\texttt{$\vdash$ methodology.tex} & \\
			\texttt{$\vdash$ results.tex} & \\
			\texttt{$\vdash$ discussion.tex} & \\
			\texttt{$\vdash$ conclusion.tex} & \\
			\texttt{$\vdash$ references.bib} & (from Zotero) \\
		\end{tabular}
	\end{block}
	\begin{itemize}
		\item Copy template folder for each new paper
		\item Modify content, keep structure
		\item Saves hours per paper
	\end{itemize}
\end{frame}
	
	% ============================================================
	% SLIDE 41: AI TOOLS FOR DRAFTING - ACCELERATE WRITING
	% ============================================================
\begin{frame}{AI Tools - Accelerate Writing to 50 Papers}
	\scriptsize
	\begin{block}{Using AI for Volume}
		\begin{itemize}
			\item \textbf{Literature review summaries:} ChatGPT summarizes 20 papers in 5 minutes
			\item \textbf{First draft generation:} Provide outline, AI writes draft
			\item \textbf{Editing and polishing:} Improve clarity and grammar
			\item \textbf{Reference formatting:} Ensure BibTeX is correct
		\end{itemize}
	\end{block}
	\begin{block}{Sample AI Prompt for Review Paper}
		\begin{center}
			\ttfamily
			Write a 3000-word literature review on \\
			"Machine Learning for Financial Stability" \\
			including sections on: (1) anomaly detection, \\
			(2) stress testing, (3) payment systems. \\
			Cite 50-80 references (provide titles).
		\end{center}
		\begin{itemize}
			\item AI generates first draft in 10-15 minutes
			\item You edit and verify (adds 2-3 hours per paper)
			\item 50 papers × 3 hours = 150 hours (doable in 2 years)
		\end{itemize}
	\end{block}
\end{frame}
	% ============================================================
	% SLIDE 42: COLLABORATION TOOLS - GITHUB
	% ============================================================
	\begin{frame}{GitHub - Collaboration for Volume}
		\scriptsize
		\begin{block}{Why GitHub for Research}
			\begin{itemize}
				\item Version control for all 50 papers
				\item Collaborate with co-authors without emailing files
				\item Public profile shows research activity
				\item Code and data can be shared
			\end{itemize}
		\end{block}
		\begin{block}{GitHub Workflow}
			\begin{enumerate}
				\item Create repository for each paper
				\item Push LaTeX code, figures, .bib file
				\item Add co-authors as collaborators
				\item Track changes (who wrote what)
				\item Merge contributions automatically
			\end{enumerate}
		\end{block}
	\end{frame}
	
	% ============================================================
	% SLIDE 43: 50 PAPERS - REALISTIC BREAKDOWN BY TYPE
	% ============================================================
	\begin{frame}{50 Papers - Realistic Breakdown by Type}
		\scriptsize
		\begin{block}{Target Mix for 50 Papers}
			\centering
			\begin{tabular}{|p{4cm}|p{3cm}|p{3cm}|}
				\hline
				\textbf{Paper Type} & \textbf{Number} & \textbf{Citations Each} \\
				\hline
				Review Papers & 12 & 150-300 \\
				\hline
				Original Research (Journals) & 15 & 30-80 \\
				\hline
				Conference Proceedings & 15 & 10-30 \\
				\hline
				Book Chapters & 5 & 10-30 \\
				\hline
				Short Communications & 3 & 5-20 \\
				\hline
				\textbf{Total} & \textbf{50} & \textbf{1000-2000+} \\
				\hline
			\end{tabular}
		\end{block}
		\begin{block}{Year-by-Year Breakdown}
			\begin{itemize}
				\item Year 1: 12 papers (3 reviews, 4 original, 4 conference, 1 chapter)
				\item Year 2: 18 papers (4 reviews, 6 original, 6 conference, 2 chapters)
				\item Year 3: 20 papers (5 reviews, 5 original, 5 conference, 2 chapters, 3 short)
			\end{itemize}
		\end{block}
	\end{frame}
	
	% ============================================================
	% SLIDE 44: TIME MANAGEMENT - 30 MINUTES DAILY FOR 50 PAPERS
	% ============================================================
	\begin{frame}{Time Management - 30 Minutes Daily for 50 Papers}
		\scriptsize
		\begin{block}{Daily Habits (30 minutes)}
			\begin{itemize}
				\item \textbf{Minutes 0-5:} Check email for review invitations, accept/decline
				\item \textbf{Minutes 5-10:} Check Google Scholar for new citations
				\item \textbf{Minutes 10-20:} Write 300-500 words on current paper
				\item \textbf{Minutes 20-25:} Save 5-10 new references to Zotero
				\item \textbf{Minutes 25-30:} Post 1 update on LinkedIn or share a paper
			\end{itemize}
		\end{block}
		\begin{block}{Weekly (2-3 hours total)}
			\begin{itemize}
				\item Submit 1 paper every 2-3 weeks
				\item Complete 1 peer review per week (2-3 hours)
				\item Review collaborator contributions and provide feedback
			\end{itemize}
		\end{block}
		\begin{block}{Monthly (8-10 hours)}
			\begin{itemize}
				\item Submit 2-3 papers
				\item Complete 4-5 peer reviews
				\item Update citation dashboard
			\end{itemize}
		\end{block}
	\end{frame}
	
	% ============================================================
	% SLIDE 45: BUDGET ALLOCATION FOR 50 PAPERS
	% ============================================================
	\begin{frame}{Budget Allocation - \$15-25K for 50 Papers}
		\scriptsize
		\begin{block}{Where the Money Goes}
			\centering
			\begin{tabular}{|p{4cm}|p{3cm}|p{3cm}|}
				\hline
				\textbf{Expense} & \textbf{Low Budget} & \textbf{High Budget} \\
				\hline
				Journal APCs (25 papers) & \$5,000 (\$200 avg) & \$12,500 (\$500 avg) \\
				\hline
				Conference registrations (15) & \$4,500 (\$300 ea) & \$9,000 (\$600 ea) \\
				\hline
				AI tools (36 months) & \$500 & \$1,500 \\
				\hline
				Professional editing (optional) & \$1,000 & \$3,000 \\
				\hline
				Miscellaneous & \$0 & \$1,000 \\
				\hline
				\textbf{Total} & \textbf{\$11,000} & \textbf{\$27,000} \\
				\hline
			\end{tabular}
		\end{block}
		\begin{block}{How to Stay Under \$20,000}
			\begin{itemize}
				\item Prioritize no-APC subscription journals (4-8 papers at \$0)
				\item Choose conferences with lower registration (\$300-\$400)
				\item Skip professional editing (use AI tools + self-edit)
				\item Apply for fee waivers (many journals offer for developing countries)
			\end{itemize}
		\end{block}
	\end{frame}
	
	
	
	
	% ============================================================
	% SLIDE 32.6: WHY SCOPUS INDEXING IS OVERRATED
	% ============================================================
	\begin{frame}{Why Scopus Indexing Is Overrated for EB2-NIW}
		\scriptsize
		\begin{block}{The Scopus Myth}
			\begin{itemize}
				\item Many researchers think: "Scopus = automatic credibility"
				\item \textcolor{red}{WRONG. USCIS looks at CITATIONS, not database names}
				\item Scopus is just ONE database (Elsevier's product)
				\item Your paper can have 500 citations but NOT be in Scopus
			\end{itemize}
		\end{block}
		
		\begin{block}{The Real Cost-Benefit Analysis}
			\centering
			\begin{tabular}{|p{4cm}|p{3cm}|p{3cm}|}
				\hline
				\textbf{Journal Type} & \textbf{APC} & \textbf{Expected Citations} \\
				\hline
				Scopus (top tier) & \$2,000-\$6,000 & 0-50 \\
				\hline
				Scopus (mid tier) & \$800-\$2,000 & 5-30 \\
				\hline
				\textcolor{red}{Scopus (low tier)} & \textcolor{red}{\$500-\$1,000} & \textcolor{red}{0-10} \\
				\hline
				Non-Scopus (government focus) & \$0-\$500 & 50-200 \\
				\hline
				Conference proceedings & \$300-\$500 & 10-50 \\
				\hline
				Preprints + promotion & \$0 & 20-100 \\
				\hline
			\end{tabular}
		\end{block}
		
		\begin{block}{When Scopus IS Worth It}
			\begin{itemize}
				\item ONLY if APC < \$500
				\item ONLY if journal ALSO has strong .gov/.edu readership
				\item ONLY if you already have 50+ papers and need variety
				\item Otherwise, \textcolor{red}{spend your money elsewhere}
			\end{itemize}
		\end{block}
	\end{frame}
	
	
	% ============================================================
	% SLIDE 32.7: PUBLISHING DECISION MATRIX
	% ============================================================
	\begin{frame}{Publishing Decision Matrix - Where to Put Your Money}
		\scriptsize
		\begin{block}{The \$15-25K Budget Decision Guide}
			\centering
			\begin{tabular}{|p{3.5cm}|p{3cm}|p{3.5cm}|}
				\hline
				\textbf{Scenario} & \textbf{APC} & \textbf{Decision} \\
				\hline
				Scopus + high citations possible & < \$500 &  PUBLISH \\
				\hline
				Scopus + unknown citation potential & \$500-\$1,000 &  THINK TWICE \\
				\hline
				Scopus + low citation potential & > \$1,000 &  AVOID \\
				\hline
				Non-Scopus + government relevance & < \$500 &   PUBLISH \\
				\hline
				Non-Scopus + good .gov potential & \$0-\$300 &  BEST VALUE \\
				\hline
				Non-Scopus + no gov angle & ANY &  SKIP \\
				\hline
			\end{tabular}
		\end{block}
		
		\begin{block}{The Golden Rule}
			\begin{center}
				\textcolor{green}{\textbf{Government Impact + Low Cost > Scopus Indexing + High Cost}}
			\end{center}
			\begin{itemize}
				\item One .gov citation = 50 regular citations
				\item Conference papers = faster, cheaper, still citable
				\item Preprints = free citations starting day 1
				\item Your EB2-NIW profile needs CITATIONS, not expensive database names
			\end{itemize}
		\end{block}
	\end{frame}
	
	% ============================================================
	% SLIDE 46: 1000-2000 CITATIONS - REALISTIC TIMELINE
	% ============================================================
	\begin{frame}{1000-2000 Citations - Realistic Timeline}
		\scriptsize
		\begin{block}{Citation Growth Projection}
			\centering
			\begin{tabular}{|p{3cm}|p{3cm}|p{3cm}|p{3cm}|}
				\hline
				\textbf{Time} & \textbf{Papers} & \textbf{Citations} & \textbf{Citations/Paper} \\
				\hline
				Month 6 & 8 & 25 & 3 \\
				\hline
				Month 12 & 15 & 100 & 7 \\
				\hline
				Month 18 & 25 & 300 & 12 \\
				\hline
				Month 24 & 35 & 800 & 23 \\
				\hline
				Month 30 & 45 & 1500 & 33 \\
				\hline
				Month 36 & 50 & 2000 & 40 \\
				\hline
			\end{tabular}
		\end{block}
		\begin{block}{Why Citations Accelerate Over Time}
			\begin{itemize}
				\item Review papers from Year 1 get heavy citations in Years 2-3
				\item Each new paper cites your previous work
				\item Network of collaborators cites each other
				\item .gov and .edu citations accumulate slowly but are high value
			\end{itemize}
		\end{block}
	\end{frame}
	
	% ============================================================
	% SLIDE 47: COMMON MISTAKES - HIGH-VOLUME EDITION
	% ============================================================
	\begin{frame}{Common Mistakes - High-Volume Edition}
		\scriptsize
		\begin{block}{Mistakes That Kill Volume Strategy}
			\begin{itemize}
				\item \textcolor{red}{Paying \$2,000+ APC for each paper} (budget will explode)
				\item \textcolor{red}{Only targeting top-tier journals} (rejections waste time)
				\item \textcolor{red}{Writing every paper alone} (collaboration = more papers)
				\item \textcolor{red}{Not using preprints} (delays citations by 6-12 months)
				\item \textcolor{red}{Not tracking citations systematically} (can't prove impact)
				\item \textcolor{red}{Ignoring government alignment} (misses .gov citations)
			\end{itemize}
		\end{block}
		\begin{block}{Mistakes to Avoid - Peer Review}
			\begin{itemize}
				\item \textcolor{red}{Declining all review invitations} (misses recognition)
				\item \textcolor{red}{Not registering on journal portals} (no invitations)
			\end{itemize}
		\end{block}
	\end{frame}
	
	% ============================================================
	% SLIDE 48: FINAL CHECKLIST - 50 PAPERS / 2000 CITATIONS
	% ============================================================
	\begin{frame}{Final Checklist - 50 Papers   2000 Citations Ready}
		\scriptsize
		\begin{block}{Publications Checklist}
			\begin{itemize}
				\item $\square$ 50 total publications (journals  conferences chapters)
				\item $\square$ 12-15 review papers (high citation generators)
				\item $\square$ 30+ Scopus WoS indexed papers
				\item $\square$ All papers uploaded to preprint servers
			\end{itemize}
		\end{block}
		\begin{block}{Citations Checklist}
			\begin{itemize}
				\item $\square$ 1000-2000+ Google Scholar citations
				\item $\square$ 10-30 .gov citations documented
				\item $\square$ 20-50 .edu citations documented
				\item $\square$ Citation growth trend (upward)
			\end{itemize}
		\end{block}
		\begin{block}{Recognition Checklist}
			\begin{itemize}
				\item $\square$ 50+ peer reviews completed
				\item $\square$ Web of Science reviewer profile
				\item $\square$ 2-5 awards (conference best paper, outstanding reviewer)
				\item $\square$ 5-10 book chapters
			\end{itemize}
		\end{block}
		\begin{block}{Profile Management Checklist}
			\begin{itemize}
				\item $\square$ Public Google Scholar profile with all papers
				\item $\square$ ORCID with publications and reviews
				\item $\square$ ResearchGate with full texts
				\item $\square$ LinkedIn with publications section
			\end{itemize}
		\end{block}
	\end{frame}
	
	% ============================================================
	% SLIDE 49: SUCCESS METRICS - TRACK THESE NUMBERS
	% ============================================================
	\begin{frame}{Success Metrics - Track These Numbers Weekly}
		\scriptsize
		\begin{block}{Weekly Tracking (5 minutes)}
			\begin{itemize}
				\item \textbf{New citations this week:} \_\_\_\_
				\item \textbf{Total citations:} \_\_\_\_
				\item \textbf{New .gov citations:} \_\_\_\_
				\item \textbf{New .edu citations:} \_\_\_\_
				\item \textbf{Papers under review:} \_\_\_\_
				\item \textbf{Peer reviews completed this week:} \_\_\_\_
			\end{itemize}
		\end{block}
		\begin{block}{Monthly Tracking (15 minutes)}
			\begin{itemize}
				\item \textbf{New papers published:} \_\_\_\_
				\item \textbf{New papers submitted:} \_\_\_\_
				\item \textbf{h-index:} \_\_\_\_
				\item \textbf{i10-index:} \_\_\_\_
				\item \textbf{Budget spent this month:} \$\_\_\_\_
			\end{itemize}
		\end{block}
	\end{frame}
	
	% ============================================================
	% SLIDE 50: CONCLUSION - YOU CAN DO THIS
	% ============================================================
	\begin{frame}{Conclusion - You Can Build This Profile}
		\scriptsize
		\begin{block}{The Formula for 50 Papers \& 2000 Citations}
			\begin{enumerate}
				\item \textbf{Write review papers} (12-15) - each generates 150-300 citations
				\item \textbf{Collaborate widely} - each collaborator adds 2-5 papers
				\item \textbf{Use preprints} - citations start immediately
				\item \textbf{Target government priorities} - .gov citations are gold
				\item \textbf{Low-APC journals} - stay within \$15-25K budget
				\item \textbf{AI-assisted writing} - accelerate to 50 papers
				\item \textbf{Daily 30 minutes} - consistency beats intensity
			\end{enumerate}
		\end{block}
		\begin{block}{The Timeline}
			\begin{itemize}
				\item \textbf{Year 1:} 12-15 papers, 50-150 citations
				\item \textbf{Year 2:} 25-35 papers, 300-800 citations
				\item \textbf{Year 3:} 45-55 papers, 1000-2000+ citations
			\end{itemize}
		\end{block}
		\begin{block}{Final Word}
			\textcolor{blue}{Impact > Journal Name. Government Angle > Everything.}
			
			\textcolor{blue}{Volume creates citations. Citations create recognition.}
			
			\textcolor{green}{Start today. Thirty minutes. One paper at a time. 50 papers in 3 years.}
		\end{block}
	\end{frame}
	
\end{document}